\documentclass[11pt,a4paper]{article}

\usepackage[T1]{fontenc}
\usepackage[utf8]{inputenc}
\usepackage{newtxtext,newtxmath}
\usepackage[varqu,varl,scaled=0.95]{zi4}
\usepackage[a4paper,margin=2cm,top=2.2cm,bottom=2cm]{geometry}
\usepackage{microtype}
\usepackage{xcolor}
\usepackage{siunitx}
\usepackage{booktabs}
\usepackage{tabularx}
\usepackage{array}
\usepackage{enumitem}
\usepackage{titlesec}
\usepackage{fancyhdr}
\usepackage[american]{circuitikz}
\usepackage[hidelinks]{hyperref}

\definecolor{accent}{HTML}{1F3A5F}
\definecolor{lcdbg}{HTML}{DDE8C8}
\definecolor{note}{HTML}{F3F4F6}

\sisetup{group-separator={,},group-minimum-digits=4,per-mode=symbol}
\setlist{itemsep=2pt,topsep=3pt}
\setlength{\parindent}{0pt}
\setlength{\parskip}{5pt plus 1pt}
\renewcommand{\arraystretch}{1.2}
\renewcommand{\tabularxcolumn}[1]{m{#1}}
\newcolumntype{L}{>{\raggedright\arraybackslash}X}
\newcolumntype{D}{>{\centering\arraybackslash}m{2.55cm}}
\newcommand{\code}[1]{\texttt{#1}}
\newcommand{\lcd}[2]{{\setlength{\fboxsep}{2.5pt}\colorbox{lcdbg}{\parbox{2.25cm}{\raggedright\ttfamily\scriptsize\strut #1\\\strut #2}}}}

\titleformat{\section}{\large\sffamily\bfseries\color{accent}}{\thesection}{0.6em}{}
\titlespacing*{\section}{0pt}{12pt}{4pt}

\pagestyle{fancy}
\fancyhf{}
\lhead{\small IR Remote Cloner: Evaluation Prep}
\rhead{\small\thepage}
\renewcommand{\headrulewidth}{0.4pt}

\hypersetup{pdftitle={IR Remote Cloner: Evaluation Prep},pdfauthor={Shreyas and Yash}}

\begin{document}

{\LARGE\sffamily\bfseries\color{accent} IR Remote Cloner: Evaluation Prep\par}
\vspace{2pt}
{\small Shreyas (EC24BT033) and Yash (EC24BT041), EE302L, November 2026\par}
\vspace{6pt}

\section{Read this first}

The evaluation checks that the cloner works, not the code. So the goal is a clean live demo, and both of us should be able to run it alone. Three things can spoil it:

\begin{enumerate}
\item \textbf{Aim.} Most failed sends during development were down to aim. The LED beam is narrow, and a few degrees off at \qty{1}{\metre} loses the projector. Stay within \qty{1}{\metre}, point the LED straight at the projector's IR sensor, and hold it there until the send finishes.
\item \textbf{SW1 must be OFF.} SW1 shares P1.0 with the carrier. When it is ON, the carrier is shorted to ground and nothing is transmitted.
\item \textbf{A remote that is not NEC format.} Only NEC buttons can be saved. Our only real test was the Epson remote, so read Section~\ref{sec:other} in case the professor hands us another one.
\end{enumerate}

\colorbox{note}{\parbox{\dimexpr\linewidth-2\fboxsep}{\textbf{The 30-second explanation.} This board learns a button from an IR remote and then replaces the remote. A receiver module gives the 8051 the timing of the signal. The 8051 decodes it and saves the code in on-chip EEPROM, so it survives power-off. To send, it rebuilds the frame and drives an IR LED with a \qty{38}{\kilo\hertz} carrier from Timer~2. There are 8 slots chosen with three switches, REC and SEND buttons, and an LCD for status. It needs no computer.}}

\section{Before the evaluation}

\textbf{Bring:} Pt-51 with LCD, soldered board and its six wires, USB cable and a power source, the Epson remote with working batteries, a written slot map, and the laptop, USB-TTL converter and jumpers for the optional tests.

\textbf{The day before:}
\begin{itemize}[label=$\square$]
\item Confirm which projector and remote will be used. If it is not the Epson, test on that device first.
\item Record the buttons and fill in the slot map (Section~\ref{sec:slots}). The LCD shows codes such as \code{8355-90}, not button names, so the map is the only record of which slot is which. Leave slot 8 empty for the live recording.
\item Find the projector's IR sensor and the spot we will send from.
\item Each of us runs the whole of Section~\ref{sec:demo} once alone, including the power cycle.
\item Try any other remotes we can find, so we already know which ones give \code{NOT NEC}.
\item If we might show the laptop tests, run them now (Section~\ref{sec:laptop}). The loopback has not been run on the final firmware yet.
\item Do not reflash unless something is broken. Whether slots survive a reflash has never been tested.
\end{itemize}

\textbf{At the bench, before the professor arrives:}
\begin{itemize}[label=$\square$]
\item All six wires are firm, and SW1 is OFF.
\item Power on and press RESET. Both LCD lines show, and the saved count is what we expect.
\item Flip SW2--SW4: the LCD follows. Press a remote button at the receiver on the home screen: LED4 blinks.
\item Send SOURCE once. The projector must be on and warmed up, not cooling down.
\end{itemize}

\newpage
\section{Demo run sheet}
\label{sec:demo}

Slot numbers below assume the planned slot map. The second LCD line on the home screen is meant to end in \code{SAVED}, but the display cuts it off at 16 characters, so \code{SA} is normal. In a slot name, \code{xx} stands for that button's command byte. The whole demo runs without the laptop.

\begin{center}
\small
\begin{tabularx}{\textwidth}{@{}c L D L@{}}
\toprule
Step & Do & LCD shows & Result and what to say\\
\midrule
1 & Power on. Press RESET if the LCD is blank. & \lcd{SEL 1: 8355-90}{IR 38kHz~~7/8 SA} & Give the 30-second explanation. Line 1 is the selected slot, line 2 the carrier and saved count.\\
\addlinespace
2 & Flip SW2--SW4. & \lcd{SEL 2: 8355-xx}{IR 38kHz~~7/8 SA} & The selection updates live. The switches are only read when a button is pressed.\\
\addlinespace
3 & Select SOURCE, aim at the sensor, tap SEND. & \lcd{SENT SLOT 2}{8355-xx} & The projector changes source. Put the real remote down first so it is clear it is not in use.\\
\addlinespace
4 & Send an arrow key a few times. & \lcd{SENT SLOT 4}{8355-xx} & The menu moves each time, so repeated sends work.\\
\addlinespace
5 & Set slot 8 (ON ON ON). Press and release REC. & \lcd{SEL 8: EMPTY}{IR 38kHz~~7/8 SA}\par\vspace{2pt}\lcd{REC SLOT 8}{PRESS REMOTE...} & Live clone. The board is now waiting for a remote.\\
\addlinespace
6 & Point the remote at the receiver, shield it from the projector, press one button once. & \lcd{SAVED SLOT 8}{8355-xx} & It timed the signal, decoded it and saved it. The name is the code: address, then command.\\
\addlinespace
7 & Aim the LED, tap SEND. & \lcd{SENT SLOT 8}{8355-xx} & The projector does what that button does, so the clone works.\\
\addlinespace
8 & Unplug USB for a few seconds, plug it back in, then SEND again. & \lcd{SEL 8: 8355-xx}{IR 38kHz~~8/8 SA} & The count is unchanged and SEND still works. The codes are in EEPROM. They lasted a week between labs.\\
\addlinespace
9 & Do this last. Select POWER, aim, hold SEND for over \qty{1}{\second}, and keep aiming. & \lcd{DOUBLE PRESS}{WAIT 6 S...}\par\vspace{2pt}\lcd{SENT SLOT 1}{8355-90} & The projector asks for confirmation, and the second send \qty{6}{\second} after the first turns it off. It needs a cool-down before it will turn on again.\\
\bottomrule
\end{tabularx}
\end{center}

A few things to remember while running it:
\begin{itemize}
\item \textbf{REC overwrites} the selected slot without asking, and \textbf{holding REC for \qty{3}{\second} erases it}. Press and release promptly.
\item \textbf{To cancel} a recording, press REC again. SEND does nothing while the board is waiting for a remote.
\item \textbf{Press the remote after} the LCD says \code{PRESS REMOTE...}, once and briefly. The recording must start at the beginning of a frame.
\item \textbf{LED4} follows the receiver only while the board is idle. It does not update during recording, so a dark LED4 at that point is normal.
\end{itemize}

\newpage
\section{Slot map}
\label{sec:slots}

ON = 1 and SW2 is the most significant bit. Replace the planned buttons with what we actually record.

\begin{center}
\small
\begin{tabular}{cccccll}
\toprule
Slot & SW2 & SW3 & SW4 & & Planned button & Code on LCD\\
\midrule
1 & OFF & OFF & OFF & & POWER & \code{8355-90}\\
2 & OFF & OFF & ON & & SOURCE & \rule{1.8cm}{0.3pt}\\
3 & OFF & ON & OFF & & ENTER & \rule{1.8cm}{0.3pt}\\
4 & OFF & ON & ON & & UP & \rule{1.8cm}{0.3pt}\\
5 & ON & OFF & OFF & & DOWN & \rule{1.8cm}{0.3pt}\\
6 & ON & OFF & ON & & LEFT & \rule{1.8cm}{0.3pt}\\
7 & ON & ON & OFF & & RIGHT & \rule{1.8cm}{0.3pt}\\
8 & ON & ON & ON & & Empty, for the live demo & \\
\bottomrule
\end{tabular}
\end{center}

\section{If something goes wrong}
\label{sec:trouble}

Check aim first. It has been the most common cause by far.

\begin{center}
\small
\begin{tabularx}{\textwidth}{@{}>{\raggedright\arraybackslash}p{4.3cm} L L@{}}
\toprule
Symptom & Likely cause & Do this\\
\midrule
\code{SENT SLOT n} shows but the projector ignores it & Aim & Move within \qty{1}{\metre}, aim at the sensor, try a few angles\\
Still nothing after re-aiming & LED not firing: SW1 ON, or P1.0 wire loose & SW1 OFF, check the wire. Watch the LED through a phone camera while pressing SEND; IR shows as a brief flash. Some cameras filter IR, so check that it sees the real remote first\\
\code{SLOT n EMPTY} & Wrong slot selected & Check the switches against the slot map\\
Projector shows the power-off dialog and stops there & It needs a second press & Hold SEND for over \qty{1}{\second} and keep aiming, or tap SEND again\\
Projector will not turn back on & Still cooling down & Wait, then tap SEND once\\
\code{PRESS REMOTE...} never changes & Receiver not seeing the remote & Press REC to cancel. On the home screen, press the remote: LED4 should blink. If not, check the batteries and the receiver wiring\\
\code{NOT NEC} with the Epson remote & Pressed too early or held down, or a stray signal & Retry with one short press after the prompt, close to the receiver and away from the lamp\\
\code{NOT NEC} with another remote & Probably not NEC format & See Section~\ref{sec:other}\\
LCD blank & Program not started & Press RESET, check the LCD is seated, adjust the contrast trimmer\\
LCD line 2 missing & LCD start-up & Unplug, plug back in, press RESET\\
LCD dim when the soldered board is connected & P1.0 unwired, so $R_2$ holds the LED on & Connect P1.0 and press RESET. If it stays dim, check $R_1$ is \qty{100}{\ohm} and not bridged by solder\\
Saved count is 0 & EEPROM is on the chip, so a swapped board has no slots & Re-record the buttons\\
Garbled text on the laptop & Converter ground not connected & Wire converter GND straight to a Pt-51 GND pin; 9600 8N1\\
\bottomrule
\end{tabularx}
\end{center}

\textbf{Last resort, reflash:} move the switch to DETACH, hold BOOT, tap RESET, release BOOT, move to ATTACH. Then run \code{make flash} in \code{Project} on the Mac, or load \code{main.hex} with FLIP on a lab PC, and press RESET. Expect to re-record if the slots are lost.

\newpage
\section{If the professor brings another remote}
\label{sec:other}

Set slot 8 and press REC as usual. One of two things happens:

\begin{center}
\small
\begin{tabularx}{\textwidth}{@{}D L@{}}
\toprule
LCD shows & What it means\\
\midrule
\lcd{SAVED SLOT 8}{xxxx-xx} & It is NEC format, so run the normal demo. To control a different device, the LED must be aimed at that device.\\
\addlinespace
\lcd{NOT NEC}{NOT SAVED} & The slot is left unchanged. Say it plainly: the board received the signal, but it decodes and stores only the NEC format, and that is a known limit of the project.\\
\bottomrule
\end{tabularx}
\end{center}

As general protocol knowledge, not something we tested: LG TVs and many generic remotes use NEC, while Sony (SIRC, \qty{40}{\kilo\hertz}), Philips (RC-5, \qty{36}{\kilo\hertz}), most Samsung TVs and air conditioners usually do not.

\textbf{Fallback with the laptop:} key \code{r} records, and if the frame is not NEC it stays in RAM as a raw capture until reset. Key \code{p} replays it, and \code{q} steps the carrier through 36, 38 and \qty{40}{\kilo\hertz}. This has only been used on the Epson remote, so try it beforehand and do not promise it will work.

\section{Optional laptop tests}
\label{sec:laptop}

Show these only if asked, and only if they passed during practice. Connect the converter, then run these on the Mac, using the port name that the first command lists:

\begin{center}
\code{ls /dev/cu.usbserial*}\\
\code{screen /dev/cu.usbserial-10 9600}
\end{center}

\begin{center}
\small
\begin{tabularx}{\textwidth}{@{}l L D@{}}
\toprule
Key & Test & Expected\\
\midrule
\code{L} & List all slots with names and codes & Printed on the laptop\\
\code{f} & Measure the carrier. Needs the jumper from P1.0 to P3.4 & \lcd{CARRIER MEASURED}{37971 Hz}\\
\code{t}, then \code{l} & Self-test: load a test frame, send it and hear it on our own receiver. Aim the LED at the receiver's lens from 10--\qty{30}{\centi\metre} & \lcd{LOOPBACK HEARD}{MATCH 67 OF 67}\\
\bottomrule
\end{tabularx}
\end{center}

The 67-of-67 result is from an earlier build, and a timing constant in the send routine has changed since. The match count may differ on the final firmware, so we should know the real number before showing it.

\newpage
\section{Circuit and wiring check}

\begin{center}
\begin{circuitikz}[x=1cm,y=1cm,line width=0.5pt,every node/.style={font=\footnotesize}]
\ctikzset{bipoles/thickness=1}
\ctikzset{resistors/scale=0.75}
\ctikzset{capacitors/scale=0.8}
\ctikzset{diodes/scale=0.8}
\ctikzset{monopoles/ground/width=0.22}

\draw[dashed,gray] (-0.05,-1.0) rectangle (6.2,14.75);
\node[anchor=north west,font=\footnotesize\sffamily\bfseries,text=gray!70!black] at (-0.05,14.7) {Pt-51 board};
\draw[dashed,gray] (6.9,1.95) rectangle (15.8,14.75);
\node[anchor=north east,font=\footnotesize\sffamily\bfseries,text=gray!70!black] at (15.75,14.7) {Soldered board};

\draw[thick] (3.2,-0.3) rectangle (5.8,13.5);
\node[align=center,font=\footnotesize\bfseries] at (4.5,8.3) {AT89C5131A};
\node[align=center,font=\scriptsize] at (4.5,7.85) {24 MHz, X1 mode};

\node[anchor=east] at (5.75,12.8) {+5\,V};
\node[anchor=east] at (5.75,12.0) {GND};
\node[anchor=east] at (5.75,9.6) {P1.0 (T2 out)};
\node[anchor=east,font=\scriptsize,text=gray!60!black] at (5.75,9.2) {SW1 must be OFF};
\node[anchor=east] at (5.75,6.0) {P3.2};
\node[anchor=east] at (5.75,4.0) {P3.5};
\node[anchor=east] at (5.75,2.7) {P3.3};
\node[anchor=east] at (5.75,1.2) {P3.0 RxD};
\node[anchor=east] at (5.75,0.4) {P3.1 TxD};

\node[anchor=west] at (3.25,12.5) {P2.0--P2.7};
\node[anchor=west] at (3.25,11.6) {P0.0};
\node[anchor=west] at (3.25,10.9) {P0.1};
\node[anchor=west] at (3.25,10.2) {P0.2};
\node[anchor=west] at (3.25,6.9) {P1.1};
\node[anchor=west] at (3.25,6.0) {P1.2};
\node[anchor=west] at (3.25,5.1) {P1.3};
\node[anchor=west] at (3.25,2.8) {P1.7};

\draw[thick] (0.1,9.7) rectangle (2.1,13.2);
\node[align=center] at (0.72,11.4) {LCD\\16$\times$2};
\node[anchor=east] at (2.05,12.5) {D0--D7};
\node[anchor=east] at (2.05,11.6) {RS};
\node[anchor=east] at (2.05,10.9) {RW};
\node[anchor=east] at (2.05,10.2) {EN};
\draw[line width=1.4pt] (2.1,12.5) -- (3.2,12.5);
\draw (2.55,12.33) -- (2.75,12.67);
\node[above] at (2.65,12.6) {8};
\draw (2.1,11.6) -- (3.2,11.6);
\draw (2.1,10.9) -- (3.2,10.9);
\draw (2.1,10.2) -- (3.2,10.2);

\draw[thick] (0.1,4.5) rectangle (2.1,7.5);
\node[align=center,font=\scriptsize] at (0.72,6.0) {ON =\\pin to\\GND};
\node[anchor=east] at (2.05,6.9) {SW2};
\node[anchor=east] at (2.05,6.0) {SW3};
\node[anchor=east] at (2.05,5.1) {SW4};
\draw (2.1,6.9) -- (3.2,6.9);
\draw (2.1,6.0) -- (3.2,6.0);
\draw (2.1,5.1) -- (3.2,5.1);

\draw[thick] (0.1,2.2) rectangle (2.1,3.4);
\node[align=center,font=\scriptsize] at (1.1,2.8) {LED4\\(IR activity)};
\draw (2.1,2.8) -- (3.2,2.8);

\draw (5.8,12.8) -- (7.4,12.8) -- (7.4,13.9) -- (13.0,13.9);
\node[above] at (9.6,13.9) {+5\,V};
\draw (5.8,12.0) -- (6.55,12.0) node[ground]{};

\draw (8.0,13.9) to[short,*-] (8.0,12.3) to[R,l={$R_2$\ \qty{2.2}{\kilo\ohm}}] (8.0,10.4) -- (8.0,9.6);
\draw (5.8,9.6) to[short,-*] (8.0,9.6) to[R,l={$R_3$\ \qty{1}{\kilo\ohm}}] (10.45,9.6);
\node[npn] (Q1) at (11.2,9.6) {};
\draw (10.45,9.6) -- (Q1.B);
\node[anchor=west,align=left] at (11.45,9.6) {$Q_1$\\BC547B};
\draw (Q1.E) -- ++(0,-0.25) node[ground]{};
\draw (Q1.C |- 0,13.9) to[short,*-] (Q1.C |- 0,13.6) to[R,l={$R_1$\ \qty{100}{\ohm}}] (Q1.C |- 0,12.1) to[full led] (Q1.C |- 0,10.7) -- (Q1.C);
\node[anchor=east,align=right] at ([xshift=-0.35cm]Q1.C |- 0,11.4) {$D_1$\\IR LED};

\draw (13.0,13.9) -- (13.0,9.6) to[R,l={$R_4$\ \qty{100}{\ohm}}] (13.0,7.7) -- (13.0,7.0);
\draw (13.0,7.7) to[short,*-] (14.55,7.7) to[eC] (14.55,4.8) node[ground]{};
\node[anchor=west,align=left,font=\scriptsize] at (14.88,6.25) {$C_1$\\\qty{4.7}{\micro\farad}};
\draw[thick] (12.2,5.0) rectangle (14.2,7.0);
\node[anchor=north,font=\scriptsize] at (13.0,6.98) {2 VS};
\node[anchor=west,font=\scriptsize] at (12.23,6.0) {3 OUT};
\node[anchor=south,font=\scriptsize] at (13.0,5.02) {1 GND};
\node[anchor=south east] at (12.9,7.03) {HS0038A2};
\draw (13.0,5.0) -- (13.0,4.75) node[ground]{};
\draw (5.8,6.0) -- (12.2,6.0);

\draw (5.8,4.0) -- (8.2,4.0) to[nopb,l={SEND}] (9.8,4.0) node[ground]{};
\draw (5.8,2.7) -- (8.2,2.7) to[nopb,l={REC}] (9.8,2.7) node[ground]{};

\draw[thick] (8.4,-0.2) rectangle (11.4,1.6);
\node[anchor=west,font=\scriptsize] at (8.45,1.2) {TXD};
\node[anchor=west,font=\scriptsize] at (8.45,0.4) {RXD};
\node[align=center,font=\scriptsize] at (10.3,0.75) {USB-TTL\\converter};
\draw[-latex] (8.4,1.2) -- (5.8,1.2);
\draw[-latex] (5.8,0.4) -- (8.4,0.4);
\draw (10.3,-0.2) -- (10.3,-0.4) node[ground]{};
\node[anchor=west,align=left,font=\scriptsize] at (11.55,0.7) {USB to computer\\(development only)};
\end{circuitikz}
\end{center}

\begin{center}
\small
\begin{tabularx}{\textwidth}{@{}l l L@{}}
\toprule
Part or signal & Pin & Check\\
\midrule
Six wires to the soldered board & & \qty{5}{\volt}, GND, P1.0, P3.2, P3.3, P3.5\\
BC547B & & Flat face towards you, left to right: C, B, E\\
HS0038A2 & & Lens towards you, left to right: 1 GND, 2 VS, 3 OUT\\
IR LED & & Long leg to $R_1$, short leg to the collector\\
USB-TTL (optional) & P3.0, P3.1 & Converter TXD to P3.0, RXD to P3.1, GND to GND\\
Carrier jumper (only for \code{f}) & P1.0 to P3.4 & Remove it afterwards\\
\bottomrule
\end{tabularx}
\end{center}

\newpage
\section{Revision: how it works}

\textbf{The signal.} A remote sends bursts of \qty{38}{\kilo\hertz} IR (marks) separated by silence (spaces), and the data is in their lengths. The carrier lets a receiver tell the signal apart from sunlight and room lighting.

\textbf{Receiving.} The HS0038A2 contains a photodiode, amplifier, \qty{38}{\kilo\hertz} band-pass filter and demodulator. Its output is low during a mark and high during a space, so P3.2 sees only the timing, never the carrier.

\textbf{Recording.} After REC, the firmware waits for the first low on P3.2, then uses Timer~0 (16-bit, \qty{0.5}{\micro\second} per tick) to time every level and stores each duration in external RAM. It stops after \qty{164}{\milli\second} of silence. Fewer than 8 durations counts as a stray flash and is ignored.

\textbf{Decoding.} An NEC frame is a \qty{9}{\milli\second} mark and \qty{4.5}{\milli\second} space, then 32 bits, then a stop mark: 67 durations in total. Every bit starts with a \qty{560}{\micro\second} mark, and the space that follows gives its value: \qty{560}{\micro\second} for 0, \qty{1690}{\micro\second} for 1. The bytes, sent least significant bit first, are address, address, command and inverted command. Epson POWER is \mbox{\code{83 55 90 6F}}: address 83 55, command 90, and 6F is 90 inverted, which is the validity check.

\textbf{Storing.} Each slot takes 13 bytes: a used flag, three code bytes and a 9-byte name. Eight slots take 104 bytes of the chip's \qty{1}{\kilo\byte} EEPROM, which keeps its contents without power. The table is copied to RAM at start-up.

\textbf{Sending.} A saved button is not replayed from the recording. The frame is rebuilt from the three stored bytes, which gives exact NEC timing, free of the receiver's distortion (marks about \qty{90}{\micro\second} long, spaces about \qty{70}{\micro\second} short), and a repeat code is added to look like a real key press. Timer~2 runs continuously in clock-out mode. For a mark the firmware connects it to P1.0, for a space it disconnects it, and Timer~0 times each one.

\textbf{Driving the LED.} A port pin cannot supply \qty{36}{\milli\ampere}, so P1.0 switches a BC547B. Port~1's internal pull-up is too weak to drive the base, which is why the \qty{2.2}{\kilo\ohm} pull-up is there. When P1.0 is low, the transistor and LED are off.

\textbf{Control.} There are no interrupts. The main loop copies the receiver state to LED4, checks both buttons, updates the LCD when the switches change and handles UART keys.

\section{What we did, in order}

\begin{enumerate}
\item \textbf{Receiver:} wired the HS0038A2 and confirmed LED4 follows a remote. The A2 pin order differs from the HS0038B, so we checked the datasheet.
\item \textbf{Carrier:} used Timer~2 clock-out instead of the planned PCA. It first measured \qty{75950}{\hertz}, because the divider runs from the full \qty{24}{\mega\hertz}. Changing RCAP2 to FF62H gave \qty{37971}{\hertz}.
\item \textbf{Self-test:} built the loopback. It was unstable with the LED at the receiver's side, and gave 67 of 67 when aimed at the lens.
\item \textbf{First clone:} captured Epson POWER and replayed it, and the projector responded. Rejected stray flashes after the lamp triggered a one-pulse capture.
\item \textbf{Slots:} added NEC decoding, rebuilt frames and 8 slots in EEPROM.
\item \textbf{Standalone:} added the LCD, switches and buttons. SW1 could not be the enable because it shares P1.0, so control moved to P3.5 and P3.3. Fixed the blank line 2 by holding EN high for \qty{1}{\milli\second}, because P0 is open-drain and has a slow pull-up.
\item \textbf{Soldered:} moved to the soldered board. The LCD dimmed with P1.0 unwired, because the pull-up held the LED on, and it was fine once fully connected. Slots survived power-off and a week between labs.
\end{enumerate}

\section{Numbers to remember}

\begin{center}
\small
\begin{tabular}{@{}ll@{\hspace{1.2cm}}ll@{}}
\toprule
Crystal, mode & \qty{24}{\mega\hertz}, X1 & NEC frame & 67 durations, 32 bits\\
Timer tick & \qty{0.5}{\micro\second} & Frame period & \qty{108}{\milli\second}\\
Carrier target, measured & \qty{38}{\kilo\hertz}, \qty{37971}{\hertz} & Slots & 8, at 13 bytes each\\
Carrier reload & RCAP2 = FF62H & Storage & \qty{1}{\kilo\byte} on-chip EEPROM\\
Carrier formula & $24\,\mathrm{MHz} / (4 \times 158)$ & LED current (calculated) & \qty{36}{\milli\ampere}\\
Tested range & about \qty{1}{\metre} & LED, base, pull-up resistors & \qty{100}{\ohm}, \qty{1}{\kilo\ohm}, \qty{2.2}{\kilo\ohm}\\
Epson POWER & \code{83 55 90 6F} & Receiver filter & \qty{100}{\ohm}, \qty{4.7}{\micro\farad}\\
\bottomrule
\end{tabular}
\end{center}

\section{If the professor asks}

\small
\begin{tabularx}{\textwidth}{@{}>{\raggedright\arraybackslash}p{5cm} L@{}}
\toprule
Question & Answer\\
\midrule
Why a \qty{38}{\kilo\hertz} carrier? & So the receiver's band-pass filter can reject steady IR from sunlight and lamps.\\
Why not toggle the pin in software? & The half period is \qty{13}{\micro\second}, about 26 machine cycles, which leaves too little time to do anything else. Timer~2 does it in hardware.\\
Why a transistor? & A port pin cannot supply the LED current. The BC547B switches \qty{36}{\milli\ampere} using about \qty{1.3}{\milli\ampere} of base current.\\
Why \qty{100}{\ohm} on the LED? & It sets the current. Pins are high during reset and boot, so the LED stays on until the program starts, and \qty{47}{\ohm} would exceed a \qty{0.25}{\watt} resistor's rating.\\
Why save the code, not the recording? & It needs 3 bytes instead of 134, and the rebuilt frame has exact timing without the receiver's distortion.\\
Can it clone any remote? & No. It saves only NEC format. Others can be replayed from RAM over the laptop link, but that is untested on other brands.\\
What is the range? & About \qty{1}{\metre} when aimed. We did no formal range test. The beam is narrow; a second \qty{100}{\ohm} in parallel would roughly double the current.\\
How did you check it without a scope? & Timer~0 counted the carrier on P3.4 for \qty{131}{\milli\second}, and the loopback compared sent and received timing.\\
Are interrupts used? & No. The firmware does one job at a time and polls, which keeps the timing loops simple.\\
What language is it in? & Embedded C, compiled with SDCC, flashed over USB.\\
\bottomrule
\end{tabularx}

\end{document}
